% ============================================================
% REFERENCES
% ============================================================
\small
\begin{thebibliography}{99}
\setlength{\itemsep}{1pt}

\bibitem{aho2006}
A. V. Aho, M. S. Lam, R. Sethi, and J. D. Ullman,
\textit{Compilers: Principles, Techniques, and Tools},
2nd ed., Pearson/Addison-Wesley, 2006.

\bibitem{cooper2007}
K. D. Cooper and L. Torczon,
\textit{Engineering a Compiler},
2nd ed., Morgan Kaufmann, 2011.

\bibitem{grune2008}
D. Grune and C. J. H. Jacobs,
\textit{Parsing Techniques: A Practical Guide},
2nd ed., Springer, 2008.

\bibitem{burke1987}
M. G. Burke and G. A. Fisher,
``A practical method for LR and LL syntactic error diagnosis and recovery,''
\textit{ACM Transactions on Programming Languages and Systems}, vol. 9, no. 2, pp. 164--197, 1987.

\bibitem{nielson1999}
F. Nielson, H. R. Nielson, and C. Hankin,
\textit{Principles of Program Analysis},
Springer, 1999.

\bibitem{kennedy1994}
A. Kennedy,
``Dimension types,''
in \textit{Proc. European Symposium on Programming (ESOP)}, LNCS vol. 788, Springer, 1994, pp. 348--362.

\bibitem{enck2010}
W. Enck et al.,
``TaintDroid: An information-flow tracking system for realtime privacy monitoring on smartphones,''
in \textit{Proc. 9th USENIX Symposium on Operating Systems Design and Implementation (OSDI)}, 2010, pp. 393--407.

\bibitem{hipaa}
U.S. Department of Health and Human Services,
``Standards for privacy of individually identifiable health information,''
45 CFR Parts 160 and 164 (HIPAA Privacy Rule).

\bibitem{leveson1993}
N. G. Leveson and C. S. Turner,
``An investigation of the Therac-25 accidents,''
\textit{IEEE Computer}, vol. 26, no. 7, pp. 18--41, 1993.

\bibitem{mco1999}
Mars Climate Orbiter Mishap Investigation Board,
\textit{Phase I Report},
NASA, November 1999.

\bibitem{allamanis2018}
M. Allamanis, E. T. Barr, P. Devanbu, and C. Sutton,
``A survey of machine learning for big code and naturalness,''
\textit{ACM Computing Surveys}, vol. 51, no. 4, article 81, 2018.

\bibitem{gupta2017}
R. Gupta, S. Pal, A. Kanade, and S. Shevade,
``DeepFix: Fixing common C language errors by deep learning,''
in \textit{Proc. AAAI Conference on Artificial Intelligence}, 2017, pp. 1345--1351.

\bibitem{pradel2018}
M. Pradel and K. Sen,
``DeepBugs: A learning approach to name-based bug detection,''
\textit{Proc. ACM on Programming Languages}, vol. 2, no. OOPSLA, article 147, 2018.

\bibitem{goodfellow2014}
I. Goodfellow et al.,
``Generative adversarial nets,''
in \textit{Advances in Neural Information Processing Systems 27}, 2014.

\bibitem{mirza2014}
M. Mirza and S. Osindero,
``Conditional generative adversarial nets,''
arXiv:1411.1784, 2014.

\bibitem{salimans2016}
T. Salimans, I. Goodfellow, W. Zaremba, V. Cheung, A. Radford, and X. Chen,
``Improved techniques for training GANs,''
in \textit{Advances in Neural Information Processing Systems 29}, 2016.

\bibitem{arjovsky2017}
M. Arjovsky, S. Chintala, and L. Bottou,
``Wasserstein generative adversarial networks,''
in \textit{Proc. 34th International Conference on Machine Learning (ICML)}, PMLR vol. 70, 2017, pp. 214--223.

\bibitem{gulrajani2017}
I. Gulrajani, F. Ahmed, M. Arjovsky, V. Dumoulin, and A. Courville,
``Improved training of Wasserstein GANs,''
in \textit{Advances in Neural Information Processing Systems 30}, 2017.

\bibitem{choi2017}
E. Choi, S. Biswal, B. Malin, J. Duke, W. F. Stewart, and J. Sun,
``Generating multi-label discrete patient records using generative adversarial networks,''
in \textit{Proc. Machine Learning for Healthcare Conference (MLHC)}, PMLR vol. 68, 2017, pp. 286--305.

\bibitem{xu2019}
L. Xu, M. Skoularidou, A. Cuesta-Infante, and K. Veeramachaneni,
``Modeling tabular data using conditional GAN,''
in \textit{Advances in Neural Information Processing Systems 32}, 2019.

\bibitem{chawla2002}
N. V. Chawla, K. W. Bowyer, L. O. Hall, and W. P. Kegelmeyer,
``SMOTE: Synthetic minority over-sampling technique,''
\textit{Journal of Artificial Intelligence Research}, vol. 16, pp. 321--357, 2002.

\bibitem{schlegl2017}
T. Schlegl, P. Seeb\"ock, S. M. Waldstein, U. Schmidt-Erfurth, and G. Langs,
``Unsupervised anomaly detection with generative adversarial networks to guide marker discovery,''
in \textit{Proc. Information Processing in Medical Imaging (IPMI)}, LNCS vol. 10265, Springer, 2017, pp. 146--157.

\bibitem{liu2008}
F. T. Liu, K. M. Ting, and Z.-H. Zhou,
``Isolation forest,''
in \textit{Proc. IEEE International Conference on Data Mining (ICDM)}, 2008, pp. 413--422.

\bibitem{scholkopf2001}
B. Sch\"olkopf, J. C. Platt, J. Shawe-Taylor, A. J. Smola, and R. C. Williamson,
``Estimating the support of a high-dimensional distribution,''
\textit{Neural Computation}, vol. 13, no. 7, pp. 1443--1471, 2001.

\bibitem{breiman2001}
L. Breiman,
``Random forests,''
\textit{Machine Learning}, vol. 45, no. 1, pp. 5--32, 2001.

\bibitem{pedregosa2011}
F. Pedregosa et al.,
``Scikit-learn: Machine learning in Python,''
\textit{Journal of Machine Learning Research}, vol. 12, pp. 2825--2830, 2011.

\bibitem{paszke2019}
A. Paszke et al.,
``PyTorch: An imperative style, high-performance deep learning library,''
in \textit{Advances in Neural Information Processing Systems 32}, 2019.

\bibitem{kingma2015}
D. P. Kingma and J. Ba,
``Adam: A method for stochastic optimization,''
in \textit{Proc. International Conference on Learning Representations (ICLR)}, 2015.

\bibitem{sokolova2009}
M. Sokolova and G. Lapalme,
``A systematic analysis of performance measures for classification tasks,''
\textit{Information Processing \& Management}, vol. 45, no. 4, pp. 427--437, 2009.

\bibitem{fawcett2006}
T. Fawcett,
``An introduction to ROC analysis,''
\textit{Pattern Recognition Letters}, vol. 27, no. 8, pp. 861--874, 2006.

\bibitem{muchnick1997}
S. S. Muchnick,
\textit{Advanced Compiler Design and Implementation},
Morgan Kaufmann, 1997.

\end{thebibliography}
